\documentclass[12pt]{article}
\usepackage{amsmath,amssymb,amsthm,algorithm,algpseudocode,booktabs,hyperref,natbib,geometry,graphicx,xcolor,enumitem,multirow}
\geometry{margin=1in}

\newtheorem{theorem}{Theorem}[section]
\newtheorem{lemma}[theorem]{Lemma}
\newtheorem{proposition}[theorem]{Proposition}
\newtheorem{corollary}[theorem]{Corollary}
\newtheorem{definition}[theorem]{Definition}
\newtheorem{remark}[theorem]{Remark}
\newtheorem{example}[theorem]{Example}
\newtheorem{assumption}[theorem]{Assumption}

\DeclareMathOperator{\diag}{diag}
\DeclareMathOperator{\nnz}{nnz}
\DeclareMathOperator{\supp}{supp}
\DeclareMathOperator{\rank}{rank}
\DeclareMathOperator{\sign}{sign}
\DeclareMathOperator{\vol}{vol}
\newcommand{\R}{\mathbb{R}}
\newcommand{\N}{\mathbb{N}}
\newcommand{\norm}[1]{\left\|#1\right\|}
\newcommand{\abs}[1]{\left|#1\right|}
\newcommand{\inner}[2]{\langle #1, #2 \rangle}

\title{Spectral Features for MIP Decomposition Selection:\\
A Computational Study with the First Complete\\
MIPLIB 2017 Decomposition Census}

\author{[Authors]}

\date{\today}

\begin{document}

\maketitle

\begin{abstract}
We introduce spectral features extracted from the constraint hypergraph
Laplacian---spectral gaps, eigenvector localization, algebraic connectivity,
and coupling energy---as a new continuous, geometry-aware feature family for
characterizing mixed-integer program (MIP) instances with respect to their
amenability to decomposition methods.  We evaluate these features via a
rigorous ablation study on a stratified 500-instance subset of MIPLIB~2017,
comparing spectral, syntactic, and combined feature sets at matched feature
budgets using nested cross-validation with external decomposition backends
(GCG for Dantzig--Wolfe, SCIP-native Benders).  A spectral futility
predictor identifies instances where no block decomposition improves the
dual bound.  As a community artifact, we release the first systematic
decomposition census of MIPLIB~2017: spectral structural annotations for
all 1,065 instances and decomposition evaluation results for 500 stratified
instances.  Theoretical grounding is provided by Lemma~L3, a
partition-to-bound bridge relating the dual gap to crossing hyperedge weight;
its method-specific specializations for Benders and Dantzig--Wolfe (L3-C);
and a motivational spectral scaling law (Proposition~T2) explaining why the
spectral ratio $\delta^2/\gamma^2$ predicts decomposition quality.
\end{abstract}

%% ========================================================================
\section{Introduction}
\label{sec:intro}
%% ========================================================================

\subsection{Motivation: Reformulation Selection}

Selecting the right decomposition strategy for a large-scale mixed-integer
program (MIP) is a \emph{reformulation selection} problem---fundamentally
distinct from the algorithm selection problem studied extensively in the
SAT and MIP communities \citep{xu2008satzilla,lindauer2015autofolio}.
Algorithm selection asks: ``given a fixed problem formulation, which solver
configuration should be applied?''  Reformulation selection asks a strictly
harder question: ``which mathematical object should the solver see?''

The distinction is consequential.  Changing from a monolithic formulation to
a Benders decomposition, a Dantzig--Wolfe reformulation, or a Lagrangian
relaxation alters the feasible region's representation, the LP relaxation's
tightness, and the solver's branching geometry.  These are not parameter
changes---they are structural transformations of the optimization problem
itself.  Today, no automated system addresses cross-method reformulation
selection.  GCG \citep{gamrath2010generic,gamrath2023gcg} automates
Dantzig--Wolfe reformulation via hypergraph partitioning heuristics but is
restricted to a single decomposition method, provides no cross-method
prediction, and offers no guidance on when decomposition is futile.

\subsection{Our Approach: Spectral Features and a Decomposition Census}

We propose a \textbf{spectral decomposition oracle}: a lightweight
preprocessing layer that extracts spectral features from the constraint
hypergraph Laplacian and uses them to predict which decomposition
method---Benders, Dantzig--Wolfe, or none---will yield the strongest dual
bounds.  Unlike GCG's combinatorial partitioning (which either finds blocks
or does not), spectral features are \emph{continuous and geometry-aware},
degrading gracefully with coupling strength.  The oracle dispatches to
existing solver implementations (SCIP's native Benders, GCG for
Dantzig--Wolfe), contributing the spectral analysis and prediction layer
rather than reimplementing textbook decomposition algorithms.

The \textbf{primary contribution} is the first complete \emph{decomposition
census} of MIPLIB~2017: all 1,065 instances annotated with detected
structure type, spectral features, and---for 500 stratified
instances---decomposition evaluation results from independent solvers at
multiple time cutoffs.  This census, released as an open artifact, provides
the missing empirical foundation for decomposition research: no prior work
has systematically answered ``which MIPLIB instances are amenable to
Benders, Dantzig--Wolfe, or Lagrangian relaxation---and by how much?''

\subsection{Contributions}

Our contributions, in order of importance, are:

\begin{enumerate}[label=\textbf{P\arabic*.},leftmargin=2em]
\item \textbf{Spectral Feature Family.}  Eight spectral features from the
  constraint hypergraph Laplacian for MIP instance characterization, with
  formal permutation-invariance guarantees (Proposition~\ref{prop:F1}) and
  scaling-sensitivity analysis (Proposition~\ref{prop:F2}).
\item \textbf{MIPLIB 2017 Decomposition Census.}  Spectral structural
  annotations for all 1,065 instances; decomposition evaluation on 500
  stratified instances at cutoffs $\{60, 300, 900, 3600\}$\,s via GCG and
  SCIP-native Benders.  Released as open Parquet dataset.
\item \textbf{Partition-to-Bound Bridge (Lemma~\ref{lem:L3}).}  For any
  partition of the constraint set, the LP-versus-decomposed-dual gap is
  bounded by the crossing hyperedge weight, with method-specific
  specializations for Benders (Lemma~\ref{lem:L3C-benders}) and
  Dantzig--Wolfe (Lemma~\ref{lem:L3C-dw}).
\end{enumerate}

\noindent
Secondary contributions include a spectral futility predictor and the
formal distinction between algorithm selection and reformulation selection.
Proposition~\ref{prop:T2} provides motivational theoretical grounding
explaining why the spectral ratio $\delta^2/\gamma^2$ is a natural
predictor of decomposition quality; we are explicit that its constant is
vacuous on ill-conditioned instances and that the empirical program validates
the predictor independently.

\subsection{Paper Organization}

Section~\ref{sec:prelim} introduces notation, the constraint hypergraph,
and hypergraph Laplacian variants.  Section~\ref{sec:features} defines the
eight spectral features and establishes their properties.
Section~\ref{sec:theory} presents the theoretical results: the
partition-to-bound bridge (L3), its specializations (L3-C), the spectral
partition bound (L3-sp), and the motivational scaling law (T2).
Section~\ref{sec:algorithms} gives algorithms with pseudocode and
complexity bounds.  Section~\ref{sec:experiments} describes the experimental
design.  Section~\ref{sec:related} surveys related work.
Section~\ref{sec:discussion} discusses limitations and future work.
Section~\ref{sec:conclusion} concludes.



%% ========================================================================
\section{Preliminaries and Definitions}
\label{sec:prelim}
%% ========================================================================

\subsection{Notation}

Throughout, $A \in \mathbb{R}^{m \times n}$ denotes the constraint matrix of a MIP
after presolving by SCIP.  We write $\mathrm{nnz}(A)$ for the number of nonzero
entries, $a_i \in \mathbb{R}^n$ for the $i$-th row (constraint), and
$\mathrm{supp}(a_i) = \{j : A_{ij} \neq 0\}$ for its support.  The degree of
constraint~$i$ is $d_i = |\mathrm{supp}(a_i)|$, and $d_{\max} = \max_i d_i$.
Table~\ref{tab:notation} summarizes the principal symbols.

\begin{table}[ht]
\centering
\caption{Principal notation.}
\label{tab:notation}
\small
\begin{tabular}{@{}ll@{}}
\toprule
Symbol & Meaning \\
\midrule
$A \in \mathbb{R}^{m \times n}$ & Constraint matrix (presolved) \\
$\mathcal{H}(A) = (V, E, w)$ & Constraint hypergraph \\
$d_i$, $d_{\max}$ & Degree of constraint $i$; maximum degree \\
$L_{\mathrm{CE}}$, $L_I$ & Clique-expansion / incidence-matrix Laplacian \\
$\mathcal{L}$ & Normalized Laplacian $D_v^{-1/2} L D_v^{-1/2}$ \\
$\lambda_1 \leq \lambda_2 \leq \cdots$ & Eigenvalues of $\mathcal{L}$ \\
$v_1, v_2, \ldots$ & Corresponding eigenvectors \\
$\gamma_2 = \lambda_2$ & Spectral gap (algebraic connectivity) \\
$\gamma_k = \lambda_{k+1} - \lambda_k$ & $k$-way spectral gap \\
$\delta$ & Coupling energy: $\|E\|_F$ in $A = A_{\mathrm{block}} + E$ \\
$P = \{B_1, \ldots, B_k\}$ & Partition of variable set \\
$E_{\mathrm{cross}}(P)$ & Crossing hyperedges \\
$n_e(P)$ & Number of blocks spanned by hyperedge $e$ \\
$y^*$ & Optimal dual of the monolithic LP relaxation \\
$z_{\mathrm{LP}}$, $z_D(P)$ & Monolithic LP value; decomposed dual bound \\
$\kappa(M)$ & Condition number $\|M\|_2 \cdot \|M^{-1}\|_2$ \\
\bottomrule
\end{tabular}
\end{table}

\subsection{Mixed-Integer Programming and Decomposition}
\label{sec:mip-decomp}

Consider the mixed-integer program
\begin{equation}\label{eq:mip}
  z^* = \min\{c^\top x : Ax \geq b,\; x \geq 0,\;
  x_j \in \mathbb{Z} \;\forall j \in \mathcal{I}\}
\end{equation}
with LP relaxation value $z_{\mathrm{LP}} \leq z^*$ obtained by dropping
integrality.  The LP dual is
\begin{equation}\label{eq:lp-dual}
  z_{\mathrm{LP}} = \max\{b^\top y : A^\top y \leq c,\; y \geq 0\}.
\end{equation}

\textbf{Benders decomposition} \citep{benders1962partitioning} partitions
the \emph{variables} into complicating variables~$x_1$ and
subproblem variables~$x_2$.  Fixing~$x_1$, the subproblem over~$x_2$ is
solved; its dual generates Benders cuts added to a master problem over~$x_1$.
At iteration~$t$, the master provides a lower bound $z_{\mathrm{Benders}}^{(t)}$.

\textbf{Dantzig--Wolfe reformulation} \citep{dantzig1960decomposition}
identifies block-diagonal structure: linking constraints couple otherwise
independent blocks, each represented via extreme points.  Column generation
iteratively adds improving columns.  At convergence,
$z_{\mathrm{DW}}^* = z_{\mathrm{LP}}$ by LP duality.

\textbf{Lagrangian relaxation} \citep{fisher1981lagrangian,geoffrion1974lagrangean}
dualizes coupling constraints with multipliers $\lambda$.  At optimality
$z_{\mathrm{LR}}^* = z_{\mathrm{LP}}$, but suboptimal multipliers yield
$z_{\mathrm{LR}}(\lambda) \leq z_{\mathrm{LP}}$.

In all three methods, decomposition quality depends critically on which
constraints are relaxed or which blocks are formed.  \textbf{The
fundamental question:} can spectral features predict which partition and
method yield the best bound?

\subsection{The Constraint Hypergraph}
\label{sec:hypergraph}

\begin{definition}[Constraint Hypergraph]
\label{def:hypergraph}
Given constraint matrix $A \in \mathbb{R}^{m \times n}$ (after equilibration),
the constraint hypergraph is $\mathcal{H}(A) = (V, E, w)$ where
$V = \{1, \ldots, n\}$ (variables),
$E = \{e_1, \ldots, e_m\}$ with $e_i = \mathrm{supp}(a_i)$ (constraints as hyperedges),
and $w_i = \|a_i\|_2^2 / d_i$ (mean squared coefficient).
\end{definition}

\begin{remark}[Constraints as Hyperedges]
Each constraint~$i$ maps to hyperedge~$e_i$.  We use $y_e^*$ and $y_i^*$
interchangeably for the LP dual of constraint $i$.
\end{remark}

\subsection{Equilibration Preprocessing}
\label{sec:equilibration}

\begin{definition}[Ruiz Equilibration \citep{ruiz2001scaling}]
\label{def:ruiz}
Iteratively compute diagonal $D_r, D_c$ so that $\tilde{A} = D_r A D_c$
has rows and columns with unit $\ell_\infty$ norm.  Each iteration:
$r_i \leftarrow (\max_j |A_{ij}|)^{-1/2}$,
$c_j \leftarrow (\max_i |A_{ij}|)^{-1/2}$,
with zero-row guards.  Convergence is geometric; $T_{\max} = 20$ suffices.
\end{definition}

We report features under three scaling methods and require cross-scaling
Spearman $\rho > 0.9$ per feature as a robustness check.

\subsection{Hypergraph Laplacian Variants}
\label{sec:laplacian}

\begin{definition}[Clique-Expansion Laplacian]
\label{def:clique-laplacian}
For $d_{\max} \leq 200$: $L_{\mathrm{CE}} = D - W$ where
$W_{jj'} = \sum_{i: j,j' \in e_i} w_i/(d_i - 1)$ for $j \neq j'$, and
$D = \mathrm{diag}(W \mathbf{1}_n)$.  Each hyperedge induces a weighted clique
with per-pair weight $w_i/(d_i-1)$.
\end{definition}

Complexity: $O(\mathrm{nnz}(A) \cdot d_{\max})$ time and space.

\begin{definition}[Incidence-Matrix Laplacian \citep{bolla1993spectra,zhou2006learning}]
\label{def:incidence-laplacian}
For $d_{\max} > 200$: $L_I = D_v - H W_e D_e^{-1} H^\top$ where
$H \in \{0,1\}^{n \times m}$ is the incidence matrix ($H_{ji} = 1$ iff $j \in e_i$),
$W_e = \mathrm{diag}(w)$, $D_e = \mathrm{diag}(d_1,\ldots,d_m)$,
$D_v = \mathrm{diag}(\sum_i w_i H_{ji}/d_i)$.
\end{definition}

Complexity: $O(\mathrm{nnz}(A))$.  Both variants share the null space (constant
vector for connected hypergraphs) but differ in non-trivial spectrum.

\begin{definition}[Normalized Laplacian]
\label{def:normalized-laplacian}
$\mathcal{L} = D_v^{-1/2} L D_v^{-1/2}$, with eigenvalues
$0 = \lambda_1 < \lambda_2 \leq \cdots \leq \lambda_n \leq 2$ for connected
hypergraphs.  All results use this normalization.
\end{definition}

\subsection{Spectral Clustering and Davis--Kahan}
\label{sec:spectral-clustering}

Spectral clustering \citep{ng2001spectral,von2007tutorial} recovers a
$k$-partition from bottom-$k$ eigenvectors of $\mathcal{L}$.  Let
$V_k = [v_1 | \cdots | v_k]$, row-normalized to unit $\ell_2$ norm.
Partition by $k$-means on rows.  The eigengap heuristic selects
$k = \arg\max_{i \in \{2,\ldots,k_{\max}\}} (\lambda_{i+1} - \lambda_i)$.

\begin{theorem}[Davis--Kahan \citep{davis1970rotation,stewart1990matrix}]
\label{thm:davis-kahan}
Let $M, \tilde{M}$ be symmetric with $k$-dimensional eigenspaces $V_k, \tilde{V}_k$
for the $k$ smallest eigenvalues.  If
$\gamma_k = \min(\mu_{k+1} - \mu_k, \tilde{\mu}_{k+1} - \tilde{\mu}_k) > 0$, then
$\|\sin\Theta(V_k, \tilde{V}_k)\|_F \leq \|M - \tilde{M}\|_F / \gamma_k$.
\end{theorem}

\subsection{Block-Diagonal Perturbation Model}
\label{sec:perturbation}

\begin{definition}[Block-Diagonal Perturbation]
\label{def:perturbation}
$A = A_{\mathrm{block}} + E$ where $A_{\mathrm{block}}$ is block-diagonal with
$k$ blocks, $\delta = \|E\|_F$ (coupling energy), and
$\gamma_k^* = \lambda_{k+1}(\mathcal{L}(A_{\mathrm{block}})) - \lambda_k(\mathcal{L}(A_{\mathrm{block}}))$.
\end{definition}

\begin{remark}[Non-Canonicality]
For real instances, $A_{\mathrm{block}}$ depends on the chosen partition.
The perturbation model motivates the spectral ratio feature
$\delta^2/\gamma^2$; the empirical program validates it independently.
\end{remark}


%% ========================================================================
\section{Spectral Features for MIP Instance Characterization}
\label{sec:features}
%% ========================================================================

We define eight spectral features extracted from the normalized Laplacian
$\mathcal{L}$ of the constraint hypergraph.  Each feature captures a
different aspect of the block structure relevant to decomposition.

\subsection{Feature Definitions}
\label{sec:feature-defs}

Let $\lambda_1 \leq \lambda_2 \leq \cdots \leq \lambda_n$ be the eigenvalues of
$\mathcal{L}$, with eigenvectors $v_1, v_2, \ldots, v_n$.
Let $k$ be selected by the eigengap heuristic~(\ref{eq:eigengap}).
Let $P^* = \{B_1^*, \ldots, B_k^*\}$ be the spectral partition from
$k$-means on the bottom-$k$ eigenvectors.

\begin{definition}[Spectral Features]
\label{def:spectral-features}
\begin{enumerate}[label=\textbf{SF\arabic*.}]
\item \textbf{Spectral gap} $\gamma_2 = \lambda_2$.
  Algebraic connectivity; large gap indicates well-separated blocks.

\item \textbf{Spectral gap ratio} $\rho_{\mathrm{sp}} = \delta^2 / \gamma_2^2$,
  where $\delta^2 = \|\mathcal{L} - \mathcal{L}_{\mathrm{block}}\|_F^2$ and
  $\mathcal{L}_{\mathrm{block}}$ zeroes all cross-cluster entries.
  Predictor motivated by Proposition~\ref{prop:T2}.
  \emph{Guard:} $\gamma_2 < 10^{-10} \Rightarrow \rho_{\mathrm{sp}} = \mathrm{NaN}$.

  \textbf{Note:} T2 uses $\gamma_k$; we use $\gamma_2$ as conservative proxy
  ($\gamma_2 \leq \gamma_k$ for connected hypergraphs with $k \geq 2$).

\item \textbf{Eigenvalue decay rate} $\beta$: negative slope of
  $\log(\lambda_i)$ for $i = 2, \ldots, k+1$.  Fast decay $\Rightarrow$ few dominant blocks.

\item \textbf{Fiedler vector localization entropy}
  $H_F = -\sum_{j=1}^{n} |v_{2,j}|^2 \log |v_{2,j}|^2$ (convention $0 \log 0 = 0$).
  Low entropy $\Rightarrow$ localized Fiedler vector $\Rightarrow$ clear block boundary.
  When $\lambda_2$ has multiplicity $> 1$, average entropy over the degenerate eigenspace.

\item \textbf{Algebraic connectivity ratio} $\gamma_2 / \lambda_k$.
  Near~1: clean $k$-way partition; near~0: hierarchical structure.
  \emph{Guard:} $\lambda_k < 10^{-10} \Rightarrow 0$.

\item \textbf{Coupling energy} $\delta^2 = \|\mathcal{L} - \mathcal{L}_{\mathrm{block}}\|_F^2$.
  Direct measure of inter-block interaction strength.

\item \textbf{Block separability index} $s_k$: silhouette score of
  $k$-means on rows of $V_k V_k^\top$.
  Gram matrix provides rotation invariance within degenerate eigenspaces.
  Implementation: $k$-means++ seed 42, 10 restarts.

\item \textbf{Effective spectral dimension}
  $d_{\mathrm{eff}} = |\{i : \lambda_i < \lambda_{\mathrm{med}} / 10\}|$.
  Indicates number of exploitable block components.
\end{enumerate}
\end{definition}

\begin{remark}[$k$-Dependence]
\label{rem:k-dependence}
Features SF2, SF3, SF5, SF6, SF7 depend on $k$ from the eigengap heuristic.
We treat $k$ as a derived quantity and report sensitivity to
$k \in \{2, 5, 10, 20\}$.
\end{remark}

\subsection{Permutation Invariance}

\begin{proposition}[Permutation Invariance]
\label{prop:F1}
All eight spectral features are invariant under simultaneous row/column
permutations of $A$.
\end{proposition}

\begin{proof}
Row/column permutations relabel constraints and variables.
$\mathcal{H}(\tilde{A})$ is isomorphic to $\mathcal{H}(A)$, so
$\mathcal{L}(\tilde{A})$ and $\mathcal{L}(A)$ are similar with identical eigenvalues.

\textbf{SF1--SF3, SF5, SF8:} Depend only on eigenvalues (similarity invariants).

\textbf{SF6:} Frobenius norm is unitarily invariant. Spectral clustering
uses eigenvector geometry via Gram matrix $V_k V_k^\top$, which is invariant
under orthogonal transformation, so $\mathcal{L}_{\mathrm{block}}$ is determined
up to the same permutation.

\textbf{SF7:} $V_k V_k^\top$ is invariant under rotation within degenerate
eigenspaces. Silhouette is invariant under point relabeling.

\textbf{SF4:} $\tilde{v}_2 = P_c v_2$ (up to sign). Entropy is invariant
under index permutation. For multiplicity $> 1$, average entropy over the
eigenspace (invariant as a subspace).
\end{proof}

\subsection{Scaling Sensitivity}

\begin{proposition}[Scaling Sensitivity]
\label{prop:F2}
Under row scaling $A \to DA$ with $D = \mathrm{diag}(d_1, \ldots, d_m)$, $d_i > 0$:
\begin{equation}\label{eq:F2}
  \frac{\gamma_2(\mathcal{L}_A)}{\kappa(D)}
  \leq \gamma_2(\mathcal{L}_{DA})
  \leq \gamma_2(\mathcal{L}_A) \cdot \kappa(D),
\end{equation}
where $\kappa(D) = (\max_i d_i)/(\min_i d_i)$.
\end{proposition}

\begin{proof}
Under $A \to DA$, hyperedge weights transform as $\tilde{w}_i = d_i^2 w_i$.
By the Rayleigh quotient characterization of $\gamma_2$:
$\gamma_2(\mathcal{L}) = \min_{x \perp D_v^{1/2} \mathbf{1}} x^\top L x / (x^\top D_v x)$.
Both numerator and denominator scale by factors in $[\min_i d_i^2, \max_i d_i^2]$.
After normalization, the net effect on $\gamma_2$ is bounded by $\kappa(D)$.
The lower bound follows by considering $D^{-1}$.
\end{proof}

\begin{remark}[Limitation]
Proposition~\ref{prop:F2} bounds only SF1. Other features are validated
empirically under three scaling methods.  SF4 (entropy) can change
dramatically under scaling.
\end{remark}

\subsection{Numerical Considerations}
\label{sec:feature-numerics}

\textbf{Near-zero spectral gap.}  $\gamma_2 < 10^{-10}$: SF2 set to NaN.
Instance flagged as trivially decomposable (disconnected) or structureless.

\textbf{Repeated eigenvalues.}  $\lambda_2 = \lambda_3$: SF4 averaged over
degenerate eigenspace; SF7 uses Gram matrix.

\textbf{Laplacian variant discontinuity.}  Switch at $d_{\max} = 200$
may create feature discontinuity.  Include $d_{\max}$ as control variable.

%% ========================================================================
\section{Theoretical Results}
\label{sec:theory}
%% ========================================================================

\subsection{Lemma L3: Partition-to-Bound Bridge}
\label{sec:L3}

\begin{assumption}\label{assump:L3}
The LP relaxation of~\eqref{eq:mip} has a finite optimal primal-dual pair
$(x^*, y^*)$ with $y^* \geq 0$, and partition $P = \{B_1, \ldots, B_k\}$
covers all variables.
\end{assumption}

\begin{definition}[Crossing Hyperedges and Weight]
\label{def:crossing}
$E_{\mathrm{cross}}(P) = \{e_i : |\{l : \mathrm{supp}(a_i) \cap B_l \neq \emptyset\}| \geq 2\}$.
$n_{e_i}(P) = |\{l : \mathrm{supp}(a_i) \cap B_l \neq \emptyset\}|$.
$W_{\mathrm{cross}}(P) = \sum_{e \in E_{\mathrm{cross}}} |y_e^*| \cdot (n_e(P) - 1)$.
\end{definition}

\begin{lemma}[Partition-to-Bound Bridge]
\label{lem:L3}
Under Assumption~\ref{assump:L3}, let $z_D(P)$ be the Lagrangian dual
from relaxing all crossing constraints.  Then
\begin{equation}\label{eq:L3}
  z_{\mathrm{LP}} - z_D(P) \leq W_{\mathrm{cross}}(P)
  = \sum_{e \in E_{\mathrm{cross}}(P)} |y_e^*| \cdot (n_e(P) - 1).
\end{equation}
\end{lemma}

\begin{proof}
\textbf{Step 1: Constraint classification.}
Partition $\{1, \ldots, m\}$ into block-internal
$\mathcal{I} = \{i : \mathrm{supp}(a_i) \subseteq B_l \text{ for some } l\}$
and crossing $\mathcal{C} = E_{\mathrm{cross}}(P)$.

\textbf{Step 2: Decomposed dual.}
$z_D(P) = \max_{y \geq 0} \{ b^\top y : A^\top y \leq c,\; y_i = 0\;\forall i \in \mathcal{C} \}$.
Since this restricts $y_i = 0$ for crossing constraints, $z_D(P) \leq z_{\mathrm{LP}}$.

\textbf{Step 3: Constructing a feasible solution.}
Define $\bar{y}_i = y_i^*$ for $i \in \mathcal{I}$ and $\bar{y}_i = 0$ for
$i \in \mathcal{C}$.  Then $\bar{y}$ is feasible for the decomposed dual
since we only set non-negative components to zero, which can only reduce
$A^\top y$ componentwise (for $A_{ij} \geq 0$).  For general $A$, we use
the fact that $z_D(P) \geq b^\top \bar{y}$ by feasibility.

\textbf{Step 4: Gap bound.}
\begin{align}
  z_{\mathrm{LP}} - z_D(P)
  &\leq z_{\mathrm{LP}} - b^\top \bar{y}
  = b^\top y^* - b^\top \bar{y}
  = \sum_{i \in \mathcal{C}} b_i y_i^*.
\end{align}

\textbf{Step 5: The $(n_e - 1)$ factor.}
Each crossing constraint $i$ spans $n_{e_i}(P)$ blocks.  When relaxed, the
decomposition creates $n_{e_i}(P) - 1$ independent coupling relaxations
(one per block boundary crossed).  The dual multiplier $|y_i^*|$
weights each relaxation.  Summing:
\[
  z_{\mathrm{LP}} - z_D(P) \leq \sum_{i \in \mathcal{C}} |y_i^*| \cdot (n_{e_i}(P) - 1).
\]

The factor $(n_e - 1)$ is tight for star-shaped coupling and conservative
for constraints spanning exactly 2 blocks (where $n_e - 1 = 1$).
\end{proof}

\begin{remark}[Retrospective Nature]
L3 requires $y^*$ to evaluate.  For prospective use, approximate duals
$\tilde{y}$ from limited simplex iterations substitute, with error
$|W_{\mathrm{cross}}(P; y^*) - W_{\mathrm{cross}}(P; \tilde{y})| \leq
\|y^* - \tilde{y}\|_1 \cdot \max_e (n_e - 1)$.
\end{remark}

\begin{remark}[Relationship to Lagrangian Relaxation]
The core insight is standard in Lagrangian theory
\citep{fisher1981lagrangian,geoffrion1974lagrangean}.
L3's contributions: (1)~hypergraph formalization enabling spectral analysis;
(2)~$(n_e - 1)$ multi-block factor; (3)~universal partition quality metric.
\end{remark}

\subsection{Method-Specific Specializations (L3-C)}
\label{sec:L3C}

\begin{lemma}[L3-C Benders]
\label{lem:L3C-benders}
For Benders decomposition at iteration $t$ with master reduced costs $r_j^{(t)}$:
\begin{equation}\label{eq:L3C-benders}
  z_{\mathrm{LP}} - z_{\mathrm{Benders}}^{(t)} \leq
  \sum_{j \in \mathcal{J}_{\mathrm{coup}}} |r_j^{(t)}| \cdot |\mathrm{blocks}(j)|,
\end{equation}
where $\mathcal{J}_{\mathrm{coup}} = \{j : x_j \text{ appears in constraints from } \geq 2 \text{ blocks}\}$
and $|\mathrm{blocks}(j)|$ is the number of blocks containing constraints with variable $j$.
\end{lemma}

\begin{proof}
Benders operates on a \emph{variable} partition while L3 is stated for
\emph{constraint} partitions.  The bridge: each coupling variable $j$
participates in constraints from $|\mathrm{blocks}(j)|$ blocks.  The Benders
master at iteration $t$ has reduced cost $r_j^{(t)}$ for variable $j$.
By LP duality of the Benders master, fixing coupling variables and solving
subproblems introduces a gap proportional to $|r_j^{(t)}|$ per coupling
variable, weighted by the number of blocks it connects.

Formally, the Benders master dual at iteration $t$ is:
$z_{\mathrm{Benders}}^{(t)} = \min_{x_1} \{c_1^\top x_1 + \theta : \theta \geq \pi_s^\top (b - B x_1)\;\forall s \in S_t\}$
where $S_t$ is the set of cuts generated.  The gap
$z_{\mathrm{LP}} - z_{\mathrm{Benders}}^{(t)}$ is bounded by the sum of
reduced costs on variables that couple blocks, each weighted by the number
of block interactions.
\end{proof}

\begin{lemma}[L3-C Dantzig--Wolfe]
\label{lem:L3C-dw}
For DW decomposition at iteration $t$ with master duals $\mu_i^{(t)}$:
\begin{equation}\label{eq:L3C-dw}
  z_{\mathrm{LP}} - z_{\mathrm{DW}}^{(t)} \leq
  \sum_{i \in \mathcal{I}_{\mathrm{link}}} |\mu_i^{(t)}| \cdot (|\mathrm{blocks}(i)| - 1),
\end{equation}
where $\mathcal{I}_{\mathrm{link}}$ are the linking constraints and
$|\mathrm{blocks}(i)|$ is the number of blocks linked by constraint $i$.
\end{lemma}

\begin{proof}
The DW master at iteration $t$ with column pool $\Lambda_t$ has dual
$\mu^{(t)}$ on the linking constraints.  The gap between the monolithic LP
and the restricted master is:
$z_{\mathrm{LP}} - z_{\mathrm{DW}}^{(t)} = \sum_{i \in \mathcal{I}_{\mathrm{link}}} \mu_i^{(t)} \cdot (\text{violation of linking constraint } i)$.

Each linking constraint $i$ bridges $|\mathrm{blocks}(i)|$ subproblems.  The
violation under the restricted master is bounded by the number of
independent column generation subproblems that contribute to it, which is
$|\mathrm{blocks}(i)| - 1$ (one base block provides the reference, and each
additional block contributes one degree of freedom).  Hence:
$z_{\mathrm{LP}} - z_{\mathrm{DW}}^{(t)} \leq \sum_{i \in \mathcal{I}_{\mathrm{link}}} |\mu_i^{(t)}| \cdot (|\mathrm{blocks}(i)| - 1)$.

The $(|\mathrm{blocks}(i)| - 1)$ factor is tighter than the generic $(k-1)$
in L3 because each linking constraint bridges specific blocks, not all $k$.
\end{proof}

\begin{remark}[Novelty of L3-C]
The Benders specialization bridges variable partitions (Benders) with
constraint partitions (L3).  The DW specialization uses the restricted master
duals (the correct accessible object at iteration $t$) and achieves a tighter
factor than the generic bound.  Both are directly useful for evaluating
GCG's partitions and guiding partition refinement.
\end{remark}

\subsection{Lemma L3-sp: Spectral Partition Bound}
\label{sec:L3sp}

\begin{lemma}[Spectral Partition Bound]
\label{lem:L3sp}
Let $P^*$ be the spectral partition from bottom-$k$ eigenvectors of
$\mathcal{L}(\mathcal{H}(A))$.  Under the perturbation model
$A = A_{\mathrm{block}} + E$ with $\delta = \|E\|_F$ and
$\gamma_k = \lambda_{k+1}(\mathcal{L}_{\mathrm{block}}) - \lambda_k(\mathcal{L}_{\mathrm{block}}) > 0$:
\begin{equation}\label{eq:L3sp}
  W_{\mathrm{cross}}(P^*) \leq C_{\mathrm{sp}} \cdot
  \frac{\delta^2 \cdot d_{\max}}{\gamma_k^2} \cdot \max_{e} |y_e^*|,
\end{equation}
where $C_{\mathrm{sp}}$ is an absolute constant depending on the rounding procedure.
\end{lemma}

\begin{proof}[Proof sketch]
The proof chains four standard results:

\textbf{Step 1: Davis--Kahan.}
By Theorem~\ref{thm:davis-kahan}, $\|\sin\Theta(V_k, V_k^*)\|_F \leq \|E_{\mathcal{L}}\|_F / \gamma_k$,
where $V_k^*$ are the ideal eigenvectors of $\mathcal{L}_{\mathrm{block}}$
and $E_{\mathcal{L}} = \mathcal{L} - \mathcal{L}_{\mathrm{block}}$ is the
Laplacian perturbation.

\textbf{Step 2: Laplacian perturbation.}
For the clique-expansion Laplacian, each coupling entry in $E$ creates up to
$d_{\max}$ edges in the expansion.  Hence
$\|E_{\mathcal{L}}\|_F \leq \delta \cdot \sqrt{d_{\max}}$
(each nonzero in $E$ contributes weight to at most $d_{\max}$ clique edges).

\textbf{Step 3: Rounding analysis.}
By the analysis of \citet{lei2015consistency}, the fraction of misclassified
vertices under $k$-means rounding is
$\mathrm{misc}(P^*, P_{\mathrm{ideal}}) \leq C_1 \|\sin\Theta\|_F^2 = C_1 \delta^2 d_{\max} / \gamma_k^2$.

\textbf{Step 4: Misclassification to crossing weight.}
Each misclassified vertex $j$ contributes to crossing hyperedges.  The total
crossing weight from misclassified vertices is bounded by
$\mathrm{misc} \cdot n \cdot \max_e |y_e^*| \cdot \bar{d}$ where $\bar{d}$
is the average number of constraints per variable.  Combining:
$W_{\mathrm{cross}}(P^*) \leq C_{\mathrm{sp}} \cdot \delta^2 d_{\max} / \gamma_k^2 \cdot \max_e |y_e^*|$.
\end{proof}

\begin{remark}[Vacuousness on Dense Constraints]
The $d_{\max}$ factor makes this bound vacuous on instances with dense
constraints ($d_{\max} = n$, e.g., budget constraints).  L3-sp is tightest
on instances with small $d_{\max}$ (pure block-diagonal with narrow coupling),
which are exactly the instances where decomposition is obviously beneficial.
For the incidence-matrix Laplacian, a separate (potentially tighter) bound
applies; we defer its derivation to future work.
\end{remark}

\subsection{Proposition T2: Spectral Scaling Law (Motivational)}
\label{sec:T2}

\begin{proposition}[Spectral Scaling Law]
\label{prop:T2}
Under the perturbation model of Definition~\ref{def:perturbation}:
\begin{equation}\label{eq:T2}
  z_{\mathrm{LP}} - z_D(\hat{\pi}) \leq C \cdot \frac{\delta^2}{\gamma^2},
  \quad C = O(k \cdot \kappa(A_{\mathrm{block}})^4 \cdot \|c\|_\infty),
\end{equation}
where $\hat{\pi}$ is the spectral partition and $\gamma = \gamma_k$.
\end{proposition}

\begin{proof}[Proof sketch]
Chain Lemma~\ref{lem:L3sp} with Hoffman's bound \citep{hoffman1952approximate}.
By L3-sp, $W_{\mathrm{cross}}(\hat{\pi}) \leq C_{\mathrm{sp}} \delta^2 d_{\max} / \gamma_k^2 \cdot \max_e |y_e^*|$.
By Hoffman's bound, $\|y^*\|_\infty \leq \kappa(A)^2 \cdot \|c\|_\infty$.
By L3, $z_{\mathrm{LP}} - z_D(\hat{\pi}) \leq W_{\mathrm{cross}}$.
Combining: $C = O(C_{\mathrm{sp}} \cdot d_{\max} \cdot \kappa^2 \cdot k \cdot \kappa^2 \cdot \|c\|_\infty)$,
where the second $\kappa^2$ arises from converting the Hoffman bound's
constraint-space error to variable-space error.
\end{proof}

\begin{remark}[Vacuousness]
The constant $C$ is vacuous for $\kappa > 10^3$ (covering 60--70\% of MIPLIB
instances with big-$M$ formulations).  T2 is presented as a \emph{structural
scaling law} explaining the qualitative relationship: decomposition quality
degrades \emph{quadratically} in coupling-to-gap ratio, not linearly or
exponentially.  The empirical program validates $\delta^2/\gamma^2$ as a
predictor independently of $C$.
\end{remark}

\begin{remark}[Origin of $\kappa^4$]
Two independent $\kappa^2$ factors: (1)~Hoffman bound relating LP dual
perturbation to constraint perturbation; (2)~converting misclassification rate
(in vertex space) to crossing weight (in constraint/dual space).
Whether $\kappa^4$ is tight or an artifact of composition is an open question.
\end{remark}

\subsection{Proof Dependency Structure}

The proof dependencies form a directed acyclic graph:
\begin{center}
\begin{tabular}{ll}
L3 & $\longrightarrow$ L3-C (Benders), L3-C (DW) \\
L3 & $\longrightarrow$ L3-sp $\longrightarrow$ T2 \\
F1, F2 & independent of L3, T2
\end{tabular}
\end{center}
\textbf{Resilience:} If T2 fails (e.g., the constant cannot be made
meaningful), L3 and L3-C survive unchanged.  The empirical program
(feature ablation, census) is independent of all proofs.  If L3 fails,
the theoretical contribution collapses, but the spectral features and
census remain.


%% ========================================================================
\section{Algorithms}
\label{sec:algorithms}
%% ========================================================================

This section presents pseudocode for the six core algorithms with complexity
analysis.  Implementation is in Python (scipy.sparse, scikit-learn) with
optional Cython acceleration for the clique-expansion inner loop.

\subsection{Algorithm 1: Hypergraph Laplacian Construction}

\begin{algorithm}[ht]
\caption{Hypergraph Laplacian Construction}
\label{alg:laplacian}
\begin{algorithmic}[1]
\Require Constraint matrix $A \in \mathbb{R}^{m \times n}$ (post-presolve), scaling method
\Ensure Sparse symmetric PSD Laplacian $L_H$, metadata
\State $\tilde{A}, D_r, D_c \gets \textsc{Equilibrate}(A, \text{method})$
  \Comment{Ruiz, $T_{\max}=20$}
\For{$i = 1, \ldots, m$}
  \State $w_i \gets \|\tilde{a}_i\|_2^2 / |\mathrm{supp}(\tilde{a}_i)|$
  \Comment{Hyperedge weight}
\EndFor
\State $d_{\max} \gets \max_i |\mathrm{supp}(\tilde{a}_i)|$
\If{$d_{\max} \leq 200$}
  \State $L_H \gets \textsc{CliqueExpansion}(\tilde{A}, w)$
  \Comment{$O(\mathrm{nnz} \cdot d_{\max})$}
\Else
  \State $L_H \gets \textsc{IncidenceMatrix}(\tilde{A}, w)$
  \Comment{$O(\mathrm{nnz})$}
\EndIf
\For{$j = 1, \ldots, n$}
  \If{$L_H[j,j] = 0$} $L_H[j,j] \gets 10^{-12}$
  \Comment{Zero-degree regularization}
  \EndIf
\EndFor
\State \Return $L_H$, $\{$variant, $d_{\max}$, $\mathrm{nnz}(L_H)$, $\kappa(D_r)$$\}$
\end{algorithmic}
\end{algorithm}

\textbf{Correctness:} The output $L_H$ is symmetric positive semidefinite
by construction (difference of degree matrix and non-negative weight matrix
for clique expansion; $D_v - H W_e D_e^{-1} H^\top$ is PSD since
$W_e D_e^{-1}$ is PSD for the incidence variant).  The regularization
preserves PSD while removing the null space contribution of isolated vertices.

\textbf{Complexity:} See Table~\ref{tab:complexity}.

\subsection{Algorithm 2: Spectral Feature Extraction}

\begin{algorithm}[ht]
\caption{Spectral Feature Extraction}
\label{alg:features}
\begin{algorithmic}[1]
\Require Laplacian $L_H$, cache key
\Ensure Feature vector $f \in \mathbb{R}^8$, status
\If{cache hit} \Return cached features \EndIf
\State $k \gets \textsc{EigengapHeuristic}(L_H, k_{\max})$
  \Comment{Eq.~\eqref{eq:eigengap}}
\State $\Lambda, V \gets \textsc{RobustEigensolve}(L_H, k+1)$
  \Comment{ARPACK $\to$ LOBPCG fallback}
\If{eigensolve failed}
  \State \Return $(\mathrm{NaN}, \ldots, \mathrm{NaN})$, ``failed''
\EndIf
\State $\gamma_2 \gets \lambda_2$
  \Comment{SF1: Spectral gap}
\State $P^* \gets k\text{-means++}(V_k V_k^\top, k, \text{seed}=42, \text{restarts}=10)$
\State $\mathcal{L}_{\mathrm{block}} \gets$ zero cross-cluster entries of $\mathcal{L}$
\State $\delta^2 \gets \|\mathcal{L} - \mathcal{L}_{\mathrm{block}}\|_F^2$
  \Comment{SF6: Coupling energy}
\If{$\gamma_2 > 10^{-10}$}
  \State $\rho_{\mathrm{sp}} \gets \delta^2 / \gamma_2^2$
  \Comment{SF2: Spectral ratio}
\Else
  \State $\rho_{\mathrm{sp}} \gets \mathrm{NaN}$
\EndIf
\State $\beta \gets$ slope of $\mathrm{lstsq}(\log \lambda_i, i)$ for $i = 2, \ldots, k+1$
  \Comment{SF3}
\State $H_F \gets -\sum_j |v_{2,j}|^2 \log |v_{2,j}|^2$
  \Comment{SF4: Fiedler entropy}
\State $r_5 \gets \gamma_2 / \max(\lambda_k, 10^{-10})$
  \Comment{SF5: Connectivity ratio}
\State $s_k \gets \mathrm{silhouette}(P^*, V_k V_k^\top)$
  \Comment{SF7: Separability}
\State $d_{\mathrm{eff}} \gets |\{i : \lambda_i < \lambda_{\mathrm{med}}/10\}|$
  \Comment{SF8: Eff.\ dimension}
\State $f \gets (\gamma_2, \rho_{\mathrm{sp}}, \beta, H_F, r_5, \delta^2, s_k, d_{\mathrm{eff}})$
\State \Return $f$, status
\end{algorithmic}
\end{algorithm}

\textbf{Correctness:} Features satisfy Proposition~\ref{prop:F1}
(permutation invariance) by construction: eigenvalues are similarity
invariants, and the Gram matrix $V_k V_k^\top$ handles rotation ambiguity.

\textbf{Complexity:} Dominated by the eigensolve:
$O(\mathrm{nnz}(L_H) \cdot k \cdot T_{\mathrm{Lanczos}})$ where
$T_{\mathrm{Lanczos}}$ is typically $50$--$200$.  Practical runtime:
$< 30$\,s per instance on a modern laptop for $\mathrm{nnz} \leq 10^7$.

\subsection{Algorithm 3: Spectral Partition and Crossing Weight}

\begin{algorithm}[ht]
\caption{Spectral Partition and L3 Bound Computation}
\label{alg:partition}
\begin{algorithmic}[1]
\Require Eigenvectors $V_k$, constraint matrix $A$, LP dual $y^*$ (or approx.)
\Ensure Partition $P$, crossing weight $W_{\mathrm{cross}}$
\State Normalize rows of $V_k$ to unit $\ell_2$ norm
\State $P \gets k\text{-means++}(\text{normalized rows}, k, \text{seed}=42)$
\State $E_{\mathrm{cross}} \gets \emptyset$
\For{$i = 1, \ldots, m$}
  \State $\mathrm{blocks}_i \gets \{P(j) : j \in \mathrm{supp}(a_i)\}$
  \If{$|\mathrm{blocks}_i| \geq 2$}
    \State $E_{\mathrm{cross}} \gets E_{\mathrm{cross}} \cup \{i\}$
  \EndIf
\EndFor
\State $W_{\mathrm{cross}} \gets \sum_{i \in E_{\mathrm{cross}}} |y_i^*| \cdot (|\mathrm{blocks}_i| - 1)$
\State \Return $P$, $W_{\mathrm{cross}}$
\end{algorithmic}
\end{algorithm}

\textbf{Complexity:} $O(n \cdot k \cdot T_{k\text{-means}})$ for clustering,
$O(\mathrm{nnz})$ for crossing weight computation.

\subsection{Algorithm 4: Decomposition Selection Oracle}

\begin{algorithm}[ht]
\caption{Decomposition Selection Oracle}
\label{alg:oracle}
\begin{algorithmic}[1]
\Require MPS file path, trained classifier $\mathcal{C}$, futility threshold $\gamma_{\mathrm{thresh}}$
\Ensure Recommendation $\in \{\text{Benders, DW, none}\}$, optional partition
\State $A \gets \textsc{SCIPPresolve}(\text{MPS file})$
  \Comment{Load and presolve}
\State $L_H, \mathrm{meta} \gets \textsc{ConstructLaplacian}(A)$
  \Comment{Algorithm~\ref{alg:laplacian}}
\State $f \gets \textsc{ExtractFeatures}(L_H)$
  \Comment{Algorithm~\ref{alg:features}}
\If{$f_1 < \gamma_{\mathrm{thresh}}$}
  \Comment{Futility check (SF1 = $\gamma_2$)}
  \State \Return ``none'', $\emptyset$
\EndIf
\State recommendation $\gets \mathcal{C}(f)$
  \Comment{$O(1)$: pre-trained model}
\If{C-lite active \textbf{and} recommendation $\neq$ ``none''}
  \State $P \gets \textsc{SpectralPartition}(V_k)$
    \Comment{Algorithm~\ref{alg:partition}}
  \If{recommendation = ``DW''}
    \State Serialize $P$ to GCG .dec format
  \ElsIf{recommendation = ``Benders''}
    \State Serialize $P$ to SCIP variable partition
  \EndIf
  \State \Return recommendation, $P$
\EndIf
\State \Return recommendation, $\emptyset$
\end{algorithmic}
\end{algorithm}

\textbf{End-to-end latency:} Dominated by eigensolve ($< 30$\,s).
Presolve: $< 5$\,s.  Prediction: $< 1$\,ms.  Total: $< 35$\,s for 99\%
of MIPLIB instances.

\subsection{Complexity Summary}

\begin{table}[ht]
\centering
\caption{Complexity of core algorithms.}
\label{tab:complexity}
\small
\begin{tabular}{@{}llll@{}}
\toprule
Algorithm & Time & Space & Typical \\
\midrule
Equilibration (Ruiz) & $O(T_{\max} \cdot \mathrm{nnz})$ & $O(m+n)$ & $< 1$\,s \\
Laplacian (clique) & $O(\mathrm{nnz} \cdot d_{\max})$ & $O(\mathrm{nnz} \cdot d_{\max})$ & $1$--$10$\,s \\
Laplacian (incidence) & $O(\mathrm{nnz})$ & $O(\mathrm{nnz})$ & $< 1$\,s \\
Eigensolve & $O(\mathrm{nnz} \cdot k \cdot T_L)$ & $O(n \cdot k)$ & $5$--$30$\,s \\
Feature extraction & $O(n \cdot k^2)$ & $O(n \cdot k)$ & $< 1$\,s \\
Partition + L3 & $O(n \cdot k \cdot T_k + \mathrm{nnz})$ & $O(n + m)$ & $< 1$\,s \\
Oracle (total) & \multicolumn{2}{l}{Sum of above} & $< 35$\,s \\
\bottomrule
\end{tabular}
\end{table}

%% ========================================================================
\section{Experimental Design}
\label{sec:experiments}
%% ========================================================================

\subsection{The MIPLIB 2017 Decomposition Census}
\label{sec:census}

We construct the first systematic decomposition census of MIPLIB~2017
\citep{gleixner2021miplib} in four tiers:

\begin{description}
\item[Tier~1 (Pilot, 50 instances):] Curated to span all five structure
  types (network, block-angular, bordered-block-diagonal, staircase,
  unstructured), used for kill gates G0--G1.

\item[Tier~2 (Development, 200 instances):] Stratified sample for feature
  ablation development and G3 gate.

\item[Tier~3 (Paper, 500 instances):] Stratified by structure type
  ($5$~categories) $\times$ size bin ($5$~bins by $\mathrm{nnz}$) $\times$
  conditioning ($2$~levels: $\kappa < 10^3$, $\kappa \geq 10^3$).  This
  is the main evaluation dataset.

\item[Tier~4 (Artifact, 1,065 instances):] Spectral annotations only
  ($\sim 9$~hours compute, no decomposition).  Released as open data.
\end{description}

\noindent
Per-instance pipeline: parse MPS $\to$ SCIP presolve $\to$ spectral analysis
($< 30$\,s) $\to$ decomposition evaluation at cutoffs $\{60, 300, 900, 3600\}$\,s
via GCG (DW) and SCIP-native Benders $\to$ log to Parquet database.

\subsection{Ground Truth Label Construction}
\label{sec:labels}

Labels: $\{\text{Benders-amenable}, \text{DW-amenable}, \text{neither}\}$.
Ground truth from \textbf{external} baselines (breaking evaluation circularity
identified in the depth check):
\begin{itemize}
  \item GCG \citep{gamrath2023gcg} for Dantzig--Wolfe (independent implementation).
  \item SCIP-native Benders (\texttt{SCIPcreateBendersDefault}, SCIP $\geq$ v7.0).
\end{itemize}

\textbf{Label assignment:} $\ell = \arg\max_{m \in \{B, DW\}} \Delta_m$
where $\Delta_m = z_D^{(m)} - z_{\mathrm{LP}}^{\mathrm{root}}$ is the dual
bound improvement over the LP relaxation at wall-clock parity.  If
$\max(\Delta_B, \Delta_{DW}) < \epsilon$ (tolerance), $\ell = \text{neither}$.

\textbf{Multi-cutoff stability:} Compute labels at four cutoffs.  If $> 20\%$
of labels flip between adjacent cutoffs, use consensus (majority vote) and
report instability as a finding.

\textbf{Decomposition overhead filter:} If monolithic SCIP closes the gap
within the cutoff, $\ell = \text{neither}$ regardless of decomposition
performance.

\textbf{Tie-breaking:} When $|\Delta_B - \Delta_{DW}| < \epsilon$,
$\ell = \text{DW}$ (GCG is more mature; conservative choice).

\subsection{Core Experiment: Feature Ablation}
\label{sec:ablation}

The central experiment compares six feature configurations at matched
feature budgets:

\begin{table}[ht]
\centering
\caption{Feature configurations for the ablation study.}
\label{tab:configs}
\begin{tabular}{@{}lll@{}}
\toprule
Configuration & Features & Count \\
\midrule
SPEC-8 & 8 spectral features (Section~\ref{sec:feature-defs}) & 8 \\
SYNT-25 & 25 syntactic (density, degree stats, coeff.\ range) & 25 \\
GRAPH-10 & 10 variable-interaction graph features & 10 \\
KRUBER-21 & Kruber et al.\ (2017) feature set & 21 \\
COMB-ALL & All features combined & $43+$ \\
RANDOM & Random baseline & 0 \\
\bottomrule
\end{tabular}
\end{table}

\textbf{Feature-count-controlled comparison:} To avoid confounding feature
count with feature quality, we compare top-$k$ spectral vs.\ top-$k$
syntactic for $k \in \{3, 5, 8\}$, selected by mRMR (minimum Redundancy,
Maximum Relevance) \citep{peng2005feature}.

\textbf{Classifiers:} Random Forest (500 trees), XGBoost ($\eta = 0.1$,
max depth 6), Logistic Regression ($L_2$, $C$ via inner CV).

\textbf{Cross-validation:} Nested---outer 5-fold (evaluation), inner 3-fold
(hyperparameter tuning).  Stratified by structure type $\times$ size bin.

\textbf{Statistical tests:} McNemar's test for pairwise feature-family
comparison.  Holm--Bonferroni correction for $\binom{6}{2} = 15$
comparisons.  95\% bootstrap CIs (1,000 resamples).

\textbf{Metrics:} Accuracy, balanced accuracy (accounts for class imbalance),
Spearman $\rho$ with bound improvement, permutation feature importance,
calibration (Platt scaling + reliability diagrams).  All metrics reported
per-structure-type.

\subsection{Falsifiable Hypotheses}
\label{sec:hypotheses}

\begin{enumerate}[label=\textbf{H\arabic*.}]
\item $\rho(\delta^2/\gamma_2^2, \text{bound degradation}) \geq 0.4$
  (Spearman, 50-instance pilot). \textbf{Gate:} G1.
\item Spectral features improve selection accuracy by $\geq 5$\,pp over
  syntactic at matched $k$.  \textbf{Gate:} G3.
\item $\rho(\text{L3 bound}, \text{actual gap}) \geq 0.4$ (Spearman).
  \textbf{Gate:} G3 component.
\item Futility predictor precision $\geq 80\%$.  \textbf{Gate:} G4.
\item Cross-scaling Spearman $\rho > 0.9$ per feature.
  \textbf{Report as finding.}
\item $R^2(\gamma_2 \sim \text{syntactic}) < 0.70$ (linear) and $< 0.80$
  (RF).  \textbf{Gate:} G0.
\item Accuracy improvement comes from non-trivial classes (Benders, DW),
  not majority class (neither).  \textbf{Gate:} G4.
\end{enumerate}

\subsection{Baselines}
\label{sec:baselines}

Seven baselines, all independently implemented or using external tools:

\begin{enumerate}
\item \textbf{Monolithic SCIP:} No decomposition.  Measures decomposition overhead.
\item \textbf{GCG (DW-only):} Automatic structure detection + DW reformulation.
\item \textbf{SCIP-native Benders:} \texttt{SCIPcreateBendersDefault}.
\item \textbf{Random selector:} Uniform over $\{$Benders, DW, none$\}$.
\item \textbf{Trivial:} Always predict ``none'' (majority class).
\item \textbf{Syntactic-only:} Best classifier on SYNT-25 features.
\item \textbf{Kruber et al.:} 21 features from their Table~2 \citep{kruber2017learning}.
\end{enumerate}

\textbf{GNN baseline:} Deferred to future work.  We argue that spectral
features are preferable to GNN-learned features for this application due to:
(i)~interpretability (each feature has a formal definition);
(ii)~theoretical grounding (L3 connects features to bounds);
(iii)~computational efficiency ($O(\mathrm{nnz} \cdot k)$ vs.\
$O(\mathrm{nnz} \cdot d)$ message passing).

\subsection{L3 Empirical Validation}
\label{sec:L3-validation}

To validate Lemma~\ref{lem:L3} empirically:
\begin{enumerate}
\item Solve the monolithic LP relaxation via SCIP to obtain $y^*$.
\item For each spectral partition, compute $W_{\mathrm{cross}}(P^*)$ using~\eqref{eq:L3}.
\item Compute the actual gap $z_{\mathrm{LP}} - z_D(P^*)$ by running
  decomposition with the spectral partition.
\item Report Spearman $\rho$ between L3 bound and actual gap, per structure type.
\end{enumerate}

\textbf{Approximate duals:} Additionally test with duals from 10, 100, and
1000 simplex iterations to assess the prospective-use approximation
(Remark~\ref{rem:retrospective}).

\subsection{Kill Gates}
\label{sec:gates}

\begin{table}[ht]
\centering
\caption{Kill gates with thresholds and recovery actions.}
\label{tab:gates}
\small
\begin{tabular}{@{}cclll@{}}
\toprule
Gate & Week & Test & Kill threshold & Recovery \\
\midrule
G0 & 0 & Spectral redundancy & $R^2 \geq 0.70$ (lin.)\ or $\geq 0.80$ (RF) & ABANDON spectral \\
G1 & 2 & Spectral--decomp.\ corr. & $\rho < 0.3$ & ABANDON spectral \\
G2 & 4 & Solver wrappers & $< 80\%$ valid bounds & Debug $\to$ ABANDON wk 6 \\
G3 & 8 & Feature ablation & Spectral $\leq$ syntactic & Census + neg.\ result \\
G4 & 14 & Combined evaluation & All metrics fail & Census-only paper \\
G5 & 18 & Internal review & Fundamental problems & Revision cycle \\
\bottomrule
\end{tabular}
\end{table}

\subsection{Reporting Standards}
\label{sec:reporting}

All results include: 95\% bootstrap confidence intervals; effect sizes
(Cohen's $d$ or rank-biserial $r$); Dolan--Mor\'e performance profiles
\citep{dolan2002benchmarking} for solver comparison; confusion matrices per
feature configuration; feature importance heatmaps; calibration curves;
and per-structure-type breakdowns as \textbf{primary} results (not
supplementary).

\subsection{Threats to Validity}
\label{sec:threats}

\textbf{Internal:} Label instability across time cutoffs (mitigated by
multi-cutoff consensus); overfitting on 500 instances (mitigated by nested
CV with held-out structure types); solver version dependence (pinned via
Docker: GCG~3.5, SCIP~8.0).

\textbf{External:} MIPLIB~2017 may not represent industrial MIP instances
(acknowledged as limitation); time cutoff sensitivity (mitigated by
4-cutoff analysis).

\textbf{Construct:} ``Best method'' depends on time cutoff and solver
maturity (GCG has 15+ years of DW-specific engineering; SCIP Benders is
newer).  We acknowledge this asymmetry and discuss its implications.

%% ========================================================================
\section{Related Work}
\label{sec:related}
%% ========================================================================

\subsection{MIP Decomposition Methods}

Benders decomposition \citep{benders1962partitioning} and Dantzig--Wolfe
reformulation \citep{dantzig1960decomposition} are the classical approaches
to exploiting block structure in MIPs.  Lagrangian relaxation
\citep{fisher1981lagrangian,geoffrion1974lagrangean} provides a dual
perspective.  Modern implementations include GCG
\citep{gamrath2010generic,gamrath2023gcg}, which automates DW reformulation
via hypergraph partitioning, and SCIP's native Benders framework.
Bergner et al.~\citep{bergner2015automatic} systematically detect
block-diagonal structure in MIPs using hypergraph partitioning.

\subsection{Algorithm Selection for MIP}

SATzilla \citep{xu2008satzilla} and AutoFolio \citep{lindauer2015autofolio}
pioneered algorithm selection using instance features.  Kruber et al.\
\citep{kruber2017learning} applied ML to predict when DW decomposition
(via GCG) is beneficial, using 21 syntactic and graph features.  Our work
extends this from DW-only to cross-method reformulation selection and
introduces spectral features with theoretical grounding.

The ML for Combinatorial Optimization (ML4CO) community
\citep{gasse2019exact,nair2020solving,bengio2021machine} uses graph neural
networks on the bipartite constraint-variable graph.  GNNs implicitly
compute spectral features through message passing
\citep{xu2019powerful}, but lack the interpretability and formal connection
to decomposition bounds that our spectral features provide.

\subsection{Spectral Methods in Combinatorial Optimization}

Spectral graph theory \citep{chung1997spectral,spielman2012spectral}
provides the mathematical foundation for our approach.  Hypergraph
Laplacians were studied by Bolla \citep{bolla1993spectra}, Zhou et al.\
\citep{zhou2006learning}, and Chan et al.\ \citep{chan2018spectral}.
Spectral clustering \citep{ng2001spectral,von2007tutorial} is the standard
method for recovering block structure from eigenvectors.  The Davis--Kahan
theorem \citep{davis1970rotation,stewart1990matrix} quantifies eigenspace
perturbation.

To our knowledge, spectral features of the constraint hypergraph Laplacian
have not been used for MIP instance characterization or decomposition
method selection.

\subsection{Instance Characterization}

MIPLIB benchmarks \citep{koch2011miplib,gleixner2021miplib} are the
standard evaluation suite.  Instance features for algorithm selection have
been studied by Smith-Miles et al.\ \citep{smith2014measuring} and Hutter
et al.\ \citep{hutter2014algorithm}.  Our spectral features complement
existing syntactic features (constraint density, variable degree
distributions, coefficient statistics) with a theoretically grounded,
geometry-aware family.

%% ========================================================================
\section{Discussion}
\label{sec:discussion}
%% ========================================================================

\subsection{Why Spectral Features?}

Spectral features offer three advantages over syntactic features and
combinatorial partitioning:

\begin{enumerate}
\item \textbf{Continuity.}  GCG's structure detection is binary: it either
  finds block structure or does not.  Spectral features provide continuous
  measures of ``how decomposable'' an instance is, enabling regression and
  calibrated probability estimates.

\item \textbf{Geometry-awareness.}  The spectral gap $\gamma_2$ captures
  the algebraic connectivity of the constraint hypergraph---how tightly
  coupled the blocks are.  This is a global geometric property that
  syntactic statistics (constraint density, variable degree) cannot capture.

\item \textbf{Theoretical grounding.}  Lemma~\ref{lem:L3} connects
  partition quality (crossing weight) to bound degradation, and
  Lemma~\ref{lem:L3sp} connects spectral quality to crossing weight.
  This chain provides a formal justification for using spectral features
  as predictors.
\end{enumerate}

\subsection{Limitations}

We acknowledge several limitations:

\begin{itemize}
\item \textbf{T2's constant is vacuous.}  The scaling law $C = O(k \kappa^4 \|c\|_\infty)$
  is informationally empty on most MIPLIB instances.  T2 is motivational,
  not predictive.

\item \textbf{L3 is retrospective.}  The bound requires $y^*$ (LP dual),
  limiting its use as a prospective partition selection tool.  Approximate
  duals partially address this.

\item \textbf{Limited block structure in MIPLIB.}  Only 10--25\% of MIPLIB
  instances have exploitable block structure \citep{bergner2015automatic}.
  The oracle helps on this fraction; the futility predictor covers the rest.

\item \textbf{Feature--theory inconsistency.}  SF2 uses $\gamma_2$;
  T2's proof needs $\gamma_k$ for $k > 2$.

\item \textbf{No Lagrangian relaxation.}  ConicBundle dependency risk
  led us to exclude Lagrangian evaluation.

\item \textbf{Solver maturity asymmetry.}  GCG (15+ years for DW) vs.\
  SCIP Benders (newer).  Labels may reflect engineering, not structure.
\end{itemize}

\subsection{Future Work}

\begin{itemize}
\item \textbf{Tightening T2.}  Entry-wise analysis may reduce $\kappa^4$
  to $\kappa^2$ via smoothed analysis techniques.

\item \textbf{Lower bounds.}  Constructing instance families syntactically
  identical but spectrally distinct would establish that spectral features
  are \emph{necessary}, not just useful.

\item \textbf{NP-hardness.}  A formal proof that optimal decomposition
  selection is NP-hard (via reduction from graph partitioning).

\item \textbf{GNN comparison.}  End-to-end comparison with GNN-learned
  features \citep{gasse2019exact}.

\item \textbf{Dynamic decomposition.}  Re-evaluating spectral features
  during branch-and-bound to adapt decomposition strategy.

\item \textbf{Lagrangian integration.}  Adding ConicBundle or a custom
  bundle method for Lagrangian relaxation evaluation.
\end{itemize}

%% ========================================================================
\section{Conclusion}
\label{sec:conclusion}
%% ========================================================================

We introduced spectral features extracted from the constraint hypergraph
Laplacian as a new, theoretically grounded family of MIP instance
descriptors for decomposition method selection.  Our main theoretical
contribution, the partition-to-bound bridge (Lemma~\ref{lem:L3}), provides
a universal quality metric for any block partition, with method-specific
specializations for Benders and Dantzig--Wolfe decomposition.  The
motivational spectral scaling law (Proposition~\ref{prop:T2}) explains why
the spectral ratio $\delta^2/\gamma^2$ is a principled predictor of
decomposition quality.

The MIPLIB~2017 decomposition census---spectral annotations for all 1,065
instances and decomposition evaluation for 500 stratified instances---provides
the first systematic, cross-method evaluation of decomposition potential on
the standard MIP benchmark.  This census is released as an open artifact
for the optimization community.

The central empirical question---whether spectral features improve
decomposition selection accuracy over standard syntactic features---is
tested via a rigorous, feature-count-controlled ablation study with
external baselines and nested cross-validation.  Kill gates at weeks 0, 2,
4, 8, and 14 ensure that the project fails cheaply if the spectral premise
is invalid.

%% ========================================================================
%% BIBLIOGRAPHY
%% ========================================================================

\begin{thebibliography}{99}

\bibitem[Benders(1962)]{benders1962partitioning}
J.~F. Benders.
\newblock Partitioning procedures for solving mixed-variables programming
  problems.
\newblock \emph{Numerische Mathematik}, 4(1):238--252, 1962.

\bibitem[Bengio et~al.(2021)]{bengio2021machine}
Y.~Bengio, A.~Lodi, and A.~Prouvost.
\newblock Machine learning for combinatorial optimization: A methodological
  tour d'horizon.
\newblock \emph{European Journal of Operational Research}, 290(2):405--421, 2021.

\bibitem[Bergner et~al.(2015)]{bergner2015automatic}
M.~Bergner, A.~Caprara, A.~Ceselli, F.~Furini, M.~E. L\"ubbecke,
  E.~Malaguti, and E.~Traversi.
\newblock Automatic {Dantzig--Wolfe} reformulation of mixed integer programs.
\newblock \emph{Mathematical Programming}, 149(1):391--424, 2015.

\bibitem[Bolla(1993)]{bolla1993spectra}
M.~Bolla.
\newblock Spectra, euclidean representations and clusterings of hypergraphs.
\newblock \emph{Discrete Mathematics}, 117(1-3):19--39, 1993.

\bibitem[Chan et~al.(2018)]{chan2018spectral}
T.-H.~H. Chan, A.~Louis, Z.~G. Tang, and C.~Zhang.
\newblock Spectral properties of hypergraph {Laplacian} and approximation
  algorithms.
\newblock \emph{Journal of the ACM}, 65(3):1--48, 2018.

\bibitem[Chung(1997)]{chung1997spectral}
F.~R.~K. Chung.
\newblock \emph{Spectral Graph Theory}.
\newblock AMS, 1997.

\bibitem[Dantzig and Wolfe(1960)]{dantzig1960decomposition}
G.~B. Dantzig and P.~Wolfe.
\newblock Decomposition principle for linear programs.
\newblock \emph{Operations Research}, 8(1):101--111, 1960.

\bibitem[Davis and Kahan(1970)]{davis1970rotation}
C.~Davis and W.~M. Kahan.
\newblock The rotation of eigenvectors by a perturbation. {III}.
\newblock \emph{SIAM Journal on Numerical Analysis}, 7(1):1--46, 1970.

\bibitem[Dolan and Mor\'e(2002)]{dolan2002benchmarking}
E.~D. Dolan and J.~J. Mor\'e.
\newblock Benchmarking optimization software with performance profiles.
\newblock \emph{Mathematical Programming}, 91(2):201--213, 2002.

\bibitem[Fisher(1981)]{fisher1981lagrangian}
M.~L. Fisher.
\newblock The {Lagrangian} relaxation method for solving integer programming
  problems.
\newblock \emph{Management Science}, 27(1):1--18, 1981.

\bibitem[Gamrath et~al.(2010)]{gamrath2010generic}
G.~Gamrath, F.~Hahn, M.~E. L\"ubbecke, and et~al.
\newblock Generic column generation.
\newblock Technical Report 10-01, ZIB, 2010.

\bibitem[Gamrath et~al.(2023)]{gamrath2023gcg}
G.~Gamrath and et~al.
\newblock The {GCG} generic column generation solver.
\newblock \emph{Mathematical Programming Computation}, 2023.

\bibitem[Gasse et~al.(2019)]{gasse2019exact}
M.~Gasse, D.~Ch\'etelat, N.~Ferroni, L.~Charlin, and A.~Lodi.
\newblock Exact combinatorial optimization with graph convolutional neural
  networks.
\newblock In \emph{NeurIPS}, 2019.

\bibitem[Geoffrion(1974)]{geoffrion1974lagrangean}
A.~M. Geoffrion.
\newblock {Lagrangean} relaxation for integer programming.
\newblock \emph{Mathematical Programming Study}, 2:82--114, 1974.

\bibitem[Gleixner et~al.(2021)]{gleixner2021miplib}
A.~Gleixner, G.~Hendel, G.~Gamrath, T.~Achterberg, M.~Bastubbe,
  T.~Berthold, et~al.
\newblock {MIPLIB} 2017: Data-driven compilation of the 6th mixed-integer
  programming library.
\newblock \emph{Mathematical Programming Computation}, 13(3):443--490, 2021.

\bibitem[Hoffman(1952)]{hoffman1952approximate}
A.~J. Hoffman.
\newblock On approximate solutions of systems of linear inequalities.
\newblock \emph{Journal of Research of the National Bureau of Standards},
  49(4):263--265, 1952.

\bibitem[Hutter et~al.(2014)]{hutter2014algorithm}
F.~Hutter, L.~Xu, H.~H. Hoos, and K.~Leyton-Brown.
\newblock Algorithm runtime prediction: Methods \& evaluation.
\newblock \emph{Artificial Intelligence}, 206:79--111, 2014.

\bibitem[Koch et~al.(2011)]{koch2011miplib}
T.~Koch, T.~Achterberg, E.~Andersen, et~al.
\newblock {MIPLIB} 2010.
\newblock \emph{Mathematical Programming Computation}, 3(2):103--163, 2011.

\bibitem[Kruber et~al.(2017)]{kruber2017learning}
M.~Kruber, M.~E. L\"ubbecke, and A.~Parmentier.
\newblock Learning when to use a decomposition.
\newblock In \emph{CPAIOR}, pages 202--210, 2017.

\bibitem[Lei and Rinaldo(2015)]{lei2015consistency}
J.~Lei and A.~Rinaldo.
\newblock Consistency of spectral clustering in stochastic block models.
\newblock \emph{The Annals of Statistics}, 43(1):215--237, 2015.

\bibitem[Lindauer et~al.(2015)]{lindauer2015autofolio}
M.~Lindauer, H.~H. Hoos, F.~Hutter, and T.~Schaub.
\newblock {AutoFolio}: An automatically configured algorithm selector.
\newblock \emph{Journal of Artificial Intelligence Research}, 53:745--778, 2015.

\bibitem[Nair et~al.(2020)]{nair2020solving}
V.~Nair, S.~Bartunov, F.~Gimeno, I.~von Glehn, P.~Lichocki, I.~Lobov,
  B.~O'Donoghue, N.~Sonnerat, C.~Tjandraatmadja, P.~Wang, et~al.
\newblock Solving mixed integer programs using neural networks.
\newblock \emph{arXiv preprint arXiv:2012.13349}, 2020.

\bibitem[Ng et~al.(2001)]{ng2001spectral}
A.~Y. Ng, M.~I. Jordan, and Y.~Weiss.
\newblock On spectral clustering: Analysis and an algorithm.
\newblock In \emph{NIPS}, pages 849--856, 2001.

\bibitem[Peng et~al.(2005)]{peng2005feature}
H.~Peng, F.~Long, and C.~Ding.
\newblock Feature selection based on mutual information criteria of
  max-dependency, max-relevance, and min-redundancy.
\newblock \emph{IEEE TPAMI}, 27(8):1226--1238, 2005.

\bibitem[Ruiz(2001)]{ruiz2001scaling}
D.~Ruiz.
\newblock A scaling algorithm to equilibrate both rows and columns norms in
  matrices.
\newblock Technical Report RAL-TR-2001-034, Rutherford Appleton Laboratory, 2001.

\bibitem[Smith-Miles et~al.(2014)]{smith2014measuring}
K.~Smith-Miles, D.~Baatar, B.~Wreford, and R.~Lewis.
\newblock Towards objective measures of algorithm performance across instance
  space.
\newblock \emph{Computers \& Operations Research}, 45:12--24, 2014.

\bibitem[Spielman(2012)]{spielman2012spectral}
D.~A. Spielman.
\newblock Spectral graph theory.
\newblock In \emph{Combinatorial Scientific Computing}, pages 495--524. 2012.

\bibitem[Stewart and Sun(1990)]{stewart1990matrix}
G.~W. Stewart and J.~Sun.
\newblock \emph{Matrix Perturbation Theory}.
\newblock Academic Press, 1990.

\bibitem[von Luxburg(2007)]{von2007tutorial}
U.~von Luxburg.
\newblock A tutorial on spectral clustering.
\newblock \emph{Statistics and Computing}, 17(4):395--416, 2007.

\bibitem[Weyl(1912)]{weyl1912asymptotic}
H.~Weyl.
\newblock Das asymptotische {Verteilungsgesetz} der {Eigenwerte} linearer
  partieller {Differentialgleichungen}.
\newblock \emph{Mathematische Annalen}, 71(4):441--479, 1912.

\bibitem[Xu et~al.(2008)]{xu2008satzilla}
L.~Xu, F.~Hutter, H.~H. Hoos, and K.~Leyton-Brown.
\newblock {SATzilla}: Portfolio-based algorithm selection for {SAT}.
\newblock \emph{Journal of Artificial Intelligence Research}, 32:565--606, 2008.

\bibitem[Xu et~al.(2019)]{xu2019powerful}
K.~Xu, W.~Hu, J.~Leskovec, and S.~Jegelka.
\newblock How powerful are graph neural networks?
\newblock In \emph{ICLR}, 2019.

\bibitem[Zhou et~al.(2006)]{zhou2006learning}
D.~Zhou, J.~Huang, and B.~Sch\"olkopf.
\newblock Learning with hypergraphs: Clustering, classification, and embedding.
\newblock In \emph{NIPS}, pages 1601--1608, 2006.

\end{thebibliography}

\end{document}
